\documentclass[10pt,a4paper]{amsart}
\usepackage[margin=19mm]{geometry}
\usepackage{amsmath,amssymb,mathtools}
\usepackage[scr=rsfs,cal=euler]{mathalfa}
\usepackage{xfrac}
\usepackage[unicode,hidelinks,bookmarks=false]{hyperref}
\newcommand{\ie}{i.e.\ }
\newcommand{\cf}{cf.\ }
\newcommand{\resp}{respectively\ }
\newcommand{\Subsection}{\subsection}
\providecommand{\nobreakdash}{\nobreak\hbox{-}}
\setlength{\emergencystretch}{2em}
\multlinegap0pt
\usepackage{amssymb}
\usepackage{bm}
\usepackage{mathtools}
\usepackage{xcolor}
\usepackage{float}
\usepackage[shortlabels]{enumitem}
\usepackage{tcolorbox}
\newtheorem{prop}{Proposition}
\newcommand{\hf}{\hspace*{0.5cm}}
\title{Graph polynomial for vertex-colored embedded graphs}
\author{Somnath Basu, Dhruv Bhasin, Siddhartha Lal, Siddhartha Patra}
\date{Source: arXiv:2204.13876v2 (CC BY 4.0)}
\begin{document}
\maketitle

\noindent\emph{Source and licence.} Graph polynomial for vertex-colored embedded graphs. Version arXiv:2204.13876v2 (CC BY 4.0), distributed by arXiv under the Creative Commons Attribution 4.0 International licence, http://creativecommons.org/licenses/by/4.0/, as recorded in the arXiv licence metadata for this exact version and shown on the abstract page https://arxiv.org/abs/2204.13876v2. Full licence notice: \texttt{licenses/arxiv-2204.13876-CC-BY-4.0.txt}.\par\smallskip

\noindent\textit{Editorial scope.} Counted unit: the article's prop prop:Gamma-app (source0.tex lines 397--409). The article's complete original proof of it is delivered with it (source0.tex lines 433--466, one \texttt{proof} environment of 34 lines, which the article places away from the statement). No mathematics is written, repaired or re-derived: the statement and the proof are the article's own lines, in the article's own notation, and no retained line is substituted or re-flowed.  No further source-local context is needed: every cross-reference of every retained line resolves inside the two delivered spans.  Nothing was omitted from any retained span, so the package carries no omission record.  No retained line contains a citation, so the wrapper carries no bibliography and no bibliography entry.  No retained line contains a figure environment, an image-inclusion command or drawing code, so no image is delivered and the package declares no asset dependency.  Editorial formatting only, in addition: an AMS \texttt{amsart} preamble that restates the article's own declarations and macros used by the retained lines, namely the environments \texttt{prop} on separate counters, together with the standard packages those lines need.  The retained environments print in this document as Proposition 1 in order of appearance, whereas the article prints the same items with the numbering of its own full document.  The complete native source of the article as distributed in the arXiv e-print for this exact version is bundled unchanged as \texttt{sources/125111/source0.tex}.
\begin{prop}[Adding a pendant edge]\label{prop:Gamma-app}
Let $(\Gamma,\sigma,\mathfrak{c})$ and let $(\Gamma_\textup{app},\tilde{\sigma}, \tilde{\mathfrak{c}})$ be given such that the vertex-coloring of $\Gamma_\textup{app}$ is an extension of the vertex-coloring of $\Gamma$, i.e., $\tilde{\mathfrak{c}}|_{V(\Gamma)}=\mathfrak{c}$. Let $e_p$ denote the edge containing the pendant vertex $v_p$ and $v_0\in V(\Gamma)$. There are three possibilities for $\tilde{\mathfrak{c}}$:\\
\hf (i) $\tilde{\mathfrak{c}}(v_p)\not\in \mathfrak{c}(V(\Gamma))$;\\
\hf (ii) $\tilde{\mathfrak{c}}(v_p)=\mathfrak{c}(v_0)$;\\
\hf (iii) $\tilde{\mathfrak{c}}(v_p)\in \mathfrak{c}(V(\Gamma))$ but $\tilde{\mathfrak{c}}(v_p)\neq \mathfrak{c}(v_0)$.\\
Then, with $N=|\mathfrak{c}(V(\Gamma))|$, the following holds
\begin{eqnarray}
  {\beta}_{(\Gamma_\textup{app},\tilde{\sigma},\tilde{\mathfrak{c}})}(x) & = & (1+x)\big({\beta}_{(\Gamma,\sigma,\mathfrak{c})}(x)+(1+x)^{N-2}\big) \label{eqn:pen-new-color}\\
  {\beta}_{(\Gamma_\textup{app},\tilde{\sigma},\tilde{\mathfrak{c}})}(x) & = & {\beta}_{(\Gamma,\sigma,\mathfrak{c})}(x) \label{eqn:pen-same-color}\\
  {\beta}_{(\Gamma_\textup{app},\tilde{\sigma},\tilde{\mathfrak{c}})}(x) & = & {\beta}_{(\Gamma,\sigma,\mathfrak{c})}(x)+(1+x)^{N-2}. \label{eqn:pen-same-non-adj-color}
\end{eqnarray}
for the cases (i), (ii) and (iii) respectively.
\end{prop}

\begin{proof}[Proof of Proposition \ref{prop:Gamma-app}]
We first derive the result for case (ii). Let $c_0=\tilde{\mathfrak{c}}(v_p)=\mathfrak{c}(v_0)$. Fix a set $S$ of $i$ colors $\mathfrak{c}(V(\Gamma))$. If $c_0\not\in S$, then there is the identity bijection between $\mathscr{F}_S(\Gamma_\textup{app},\tilde{\mathfrak{c}})$ and $\mathscr{F}_S(\Gamma,\mathfrak{c})$ The corresponding contributions to $\mathscr{D}_i$'s are the same. If $c_0\in S$, then there is a natural bijection between $\mathscr{F}_S(\Gamma,\mathfrak{c})$ and $\mathscr{F}_S(\Gamma_\textup{app},\tilde{\mathfrak{c}})$; it maps $\Gamma'$ to $\Gamma'\cup e_p$. As adding an edge (via a pendant vertex) to a graph does not change the connectivity of the complement, it follows that $\mathscr{D}_i$'s are the same. This implies \eqref{eqn:pen-same-color}.\\
\hf For case (i), note that for any pair of spaces $(Y,B)$, the connected components of $Y-B$ are in bijection with the connected components of $Y/B -[B]$, i.e., the number of connected components of $Y-B$ is the rank of $\tilde{H}_0(Y/B;\mathbb{Z})$. We also note that $(\Sigma,\tilde{\sigma}(\Gamma_\textup{app}))$ and $(\Sigma,\sigma(\Gamma))$ are good pairs. For a good pair $(X,A)$, there are isomorphisms $H_j(X,A;\mathbb{Z})\cong \tilde{H}_j(X/A;\mathbb{Z})$. As $\Gamma_\textup{app}$ deformation retracts to $\Gamma$, the long exact sequence (in homology) for the triple $\sigma(\Gamma)\subset \tilde{\sigma}(\Gamma_\textup{app})\subset \Sigma$, implies that the inclusion of pairs $\iota:(\Sigma,\sigma(\Gamma))\to (\Sigma,\tilde{\sigma}(\Gamma_\textup{app}))$ induce isomorphisms
$$\iota_\ast: H_j((\Sigma,\sigma(\Gamma);\mathbb{Z})\to H_j(\Sigma,\tilde{\sigma}(\Gamma_\textup{app});\mathbb{Z})$$
In particular, when $j=0$ we obtain 
\begin{equation}\label{eqn:Dapp}
\mathscr{D}_{N+1}(\Gamma_\textup{app},\tilde{\sigma},\tilde{\mathfrak{c}})=\mathscr{D}_N(\Gamma,\sigma,\mathfrak{c})
\end{equation}
The count $\mathscr{D}_S(\Gamma_\textup{app},\tilde{\sigma},\tilde{\mathfrak{c}})$ for $1\leq |S|=j\leq N$ comes from three mutually exclusive and exhaustive cases\\
\hf (a) $S\subset \mathfrak{c}(V(\Gamma))$, and\\
\hf (b) $c_0:=\tilde{\mathfrak{c}}(v_p)\in S$, $\mathfrak{c}(v_0)\not\in S$, and\\
\hf (c) $c_0\in S$, $\mathfrak{c}(v_0)\in S$.\\
The count for (a) is $\mathscr{D}_S(\Gamma,\sigma,\mathfrak{c})$. The count for (c) is $\mathscr{D}_{S-\{c_0\}}(\Gamma,\sigma,\mathfrak{c})$, where $\mathfrak{c}(v_0)\in S$. The count for (b) is $\mathscr{D}_{S-\{c_0\}}(\Gamma,\sigma,\mathfrak{c})+1$, where $\mathfrak{c}(v_0)\not\in S$ and the $+1$ is for the isolated vertex $v_p$ in $\mathscr{F}_S(\Gamma_\textup{app},\tilde{\sigma},\widetilde{\mathfrak{c}})$. Summing over all subsets of $\mathfrak{c}(V(\Gamma))\cup\{c_0\}$ of size $j$, we obtain
\begin{equation}\label{eqn:Djapp}
\mathscr{D}_j(\Gamma_\textup{app},\tilde{\sigma},\tilde{\mathfrak{c}})=\mathscr{D}_{j}(\Gamma,\sigma,\mathfrak{c})+\mathscr{D}_{j-1}(\Gamma,\sigma,\mathfrak{c})+{N-1\choose j-1}.
\end{equation}
Combining \eqref{eqn:Dapp}, \eqref{eqn:Djapp} with the definition of the total connected induced subgraph boundary polynomial ${\beta}$, we obtain \eqref{eqn:pen-new-color}
\begin{equation*}
{\beta}_{(\Gamma_\textup{app},\tilde{\sigma},\tilde{\mathfrak{c}})}(x)=(1+x){\beta}_{(\Gamma,\sigma,\mathfrak{c})}(x)+(1+x)^{N-1}. 
\end{equation*}
\hf For case (iii), \eqref{eqn:Dapp} still holds, i.e.,
\begin{equation*}%\label{eqn:Dapp-ver2}
\mathscr{D}_{N}(\Gamma_\textup{app},\tilde{\sigma},\tilde{\mathfrak{c}})=\mathscr{D}_N(\Gamma,\sigma,\mathfrak{c})
\end{equation*}
The count $\mathscr{D}_S(\Gamma_\textup{app},\tilde{\sigma},\tilde{\mathfrak{c}})$ for $1\leq |S|=j\leq N-1$ comes from four mutually exclusive and exhaustive cases\\
\hf (a) $c_0:=\tilde{\mathfrak{c}}(v_p),\mathfrak{c}(v_0)\in S$, and\\
\hf (b) $c_0\in S$, $\mathfrak{c}(v_0)\not\in S$, and\\
\hf (c) $c_0\not\in S$, $\mathfrak{c}(v_0)\in S$, and\\
\hf (d) $c_0\not\in S$, $\mathfrak{c}(v_0)\not\in S$.\\
In all three cases (a), (c) and (d), the count is $\mathscr{D}_S(\Gamma,\sigma,\mathfrak{c})$ while for (b) it is $\mathscr{D}_S(\Gamma,\sigma,\mathfrak{c})+1$, where the $+1$ comes from the isolated vertex $v_p$ in $\mathscr{F}_S(\Gamma_\textup{app},\tilde{\sigma},\widetilde{\mathfrak{c}})$. As there are ${N-2 \choose j-1}$ such graphs satisfying (b), on summing over $S$ of size $j$, we obtain \eqref{eqn:pen-same-non-adj-color}
\begin{equation*}
{\beta}_{(\Gamma_\textup{app},\tilde{\sigma},\tilde{\mathfrak{c}})}(x)={\beta}_{(\Gamma,\sigma,\mathfrak{c})}(x)+(1+x)^{N-2}. 
\end{equation*}
\end{proof}

\end{document}
